\documentclass[journal,compsoc]{IEEEtran}

\usepackage{amsmath,amssymb,amsfonts}
\usepackage{algorithmic}
\usepackage{graphicx}
\usepackage{textcomp}
\usepackage{xcolor}
\usepackage{booktabs}
\usepackage{multirow}
\usepackage{url}
\usepackage{hyperref}
\usepackage{cite}

\hypersetup{
    colorlinks=true,
    linkcolor=blue,
    citecolor=blue,
    urlcolor=blue
}

\begin{document}

\title{The Autoregression Fallacy in Time-Critical Cyber-Physical Systems and Algorithmic Trading: An Empirical Benchmark of Non-Autoregressive Decision Models versus Foundation Large Language Models}

\author{Richie~O.~Sumual
\thanks{Richie O. Sumual is with Advanced Agentic AI \& Distributed Systems Research, Jakarta, Indonesia (e-mail: richie.sumual@research.org).}}

\markboth{IEEE Transactions on Cybernetics / IEEE Transactions on Computational Social Systems,~Vol.~XX, No.~X, September~2026}%
{Sumual: The Autoregression Fallacy in Time-Critical Cyber-Physical Systems and Algorithmic Trading}

\IEEEtitleabstractindextext{%
\begin{abstract}
The prevailing tendency to deploy autoregressive Large Language Models (LLMs) across arbitrary decision boundaries has introduced severe architectural pathologies into time-critical cyber-physical systems (CPS) and sub-second algorithmic trading. Autoregressive decoders suffer from an intrinsic generation lag: sequential next-token synthesis enforces an $\mathcal{O}(N)$ temporal execution delay and memory bandwidth saturation, inducing catastrophic state-decision drift where physical or market conditions evolve faster than policy resolution. In this paper, we present \textbf{DecisionModelBench}, a multi-domain empirical benchmark evaluating non-autoregressive decision models against foundation generative LLMs (Sahabat-AI 8B, Qwen 2.5 7B, Gemma 2 9B, and Gemma 2 2B) across distributed dual NVIDIA Tesla T4 GPU topologies and managed cloud APIs. We formulate the Continuous State Drift integral, Latency-Induced Order Book Slippage ($P_{\text{fill}} = P_{\text{signal}}(1 \pm \gamma \sqrt{\Delta t / \tau_0})$), and Negative Alpha Decay ($\alpha(\Delta t) = \alpha_0 e^{-\lambda \Delta t}$). Across three adversarial environments---tactical air defense (C-RAM/Iron Dome multi-threat discrimination under 20-missile magazine constraints and 3.5s reload windows), sub-second order book execution (\$10,000 dummy portfolio), and continuous dynamic arcade interception (Brick Breaker ball velocity vectors)---non-autoregressive decision models (Laya 421M at 55 ms, TypeSafe JEV at 160 ms, OpenJev 0.5B at 210 ms) achieve superior performance with zero token waste. Conversely, autoregressive LLMs (2,200--2,800 ms latency) precipitate structural failure: total city collapse in air defense, -\$19.41 P\&L versus +\$40.90 (Laya) in trading, and 100\% loss-of-control in arcade kinematics. We rigorously define the microstructural boundaries separating ultra-high-frequency hardware (FPGA/C++ in microsecond regimes) from sub-second decision models, demonstrating that remote cloud decision APIs outperform local 8B LLMs due to sequential memory bus bottlenecks. Finally, we formalize a Decoupled Two-Tier Cognitive Architecture uniting sub-100ms deterministic reflex triage (System 1) with out-of-band asynchronous strategic deliberation (System 2).
\end{abstract}

\begin{IEEEkeywords}
Cyber-Physical Systems, Non-Autoregressive Models, Foundation Large Language Models, Algorithmic Trading, Air Defense Simulation, Real-Time Control, Latency Slippage, Alpha Decay, Two-Tier Architecture.
\end{IEEEkeywords}}

\maketitle
\IEEEdisplaynontitleabstractindextext
\IEEEpeerreviewmaketitle

\section{Introduction}
\IEEEPARstart{T}{he} rapid evolution of generative artificial intelligence has fostered an industry-wide heuristic: the uncritical application of autoregressive Large Language Models (LLMs) as universal controllers across complex software systems~\cite{achiam2023gpt4}. Termed colloquially as the ``LLM for everything'' paradigm, this design pattern routes structured perception, triage, routing, and real-time execution queries directly through autoregressive decoder pipelines parameterized between 7 billion and 70 billion parameters~\cite{touvron2023llama}.

While autoregressive architectures exhibit remarkable sample efficiency and semantic depth in unconstrained natural language dialogues, offline document synthesis, and code generation, their deployment as closed-loop controllers in time-critical cyber-physical systems (CPS) and quantitative algorithmic trading represents a fundamental architectural fallacy. Autoregressive text generation produces output sequences token by token:
\begin{equation}
\Pr(y_{1:N} \mid x) = \prod_{i=1}^{N} \Pr(y_i \mid y_{<i}, x)
\end{equation}
Consequently, generating a structured JSON routing object or execution directive of length $N \in [30, 250]$ tokens requires $N$ sequential memory-bound forward passes through the multi-layer Transformer decoder. Even when quantized to 4-bit precision (e.g., GGUF Q4\_K\_M) and accelerated via CUDA kernel optimizations, state-of-the-art 7B--9B parameter models executed on standard enterprise accelerators (such as NVIDIA Tesla T4) demonstrate inference latencies ranging from 2,200~ms to over 6,000~ms per decision cycle.

In deterministic control theory and financial microstructure, time is an irreversible physical coordinate. Physical systems evolve continuously according to ordinary differential equations:
\begin{equation}
\frac{d S(t)}{dt} = f(S(t), u(t), t)
\end{equation}
When an agent observes state $S_t$ at time $t$ but requires an inference duration $\Delta t_{\text{inference}} = t_{\text{act}} - t$ to emit control action $u_t$, the action is applied not to state $S_t$, but to state $S_{t + \Delta t_{\text{inference}}}$. If the rate of environmental divergence $\left\|\frac{\partial S}{\partial t}\right\|$ exceeds the control margin, the issued policy becomes non-viable, outdated, or hazardous---a phenomenon defined herein as \textbf{State-Decision Drift}.

In quantitative financial trading, execution delays induce immediate negative alpha decay and adverse selection. In tactical air defense, delays in target discrimination result in ballistic impact before the kinetic interceptor departs the launch canister. In dynamic visual tracking, decision lag induces control starvation and complete kinetic failure.

To rigorously evaluate this dichotomy, this paper presents \textbf{DecisionModelBench}, a reproducible benchmark platform designed to compare non-autoregressive decision models (System~1) against foundation generative LLMs (System~2). The primary contributions of this work are as follows:
\begin{enumerate}
    \item \textbf{Mathematical Formulation of Temporal Pathology:} We provide analytical derivations for Continuous State Drift, Latency-Induced Slippage under order book diffusion, and Competitive Alpha Decay, establishing the theoretical boundaries where autoregressive inference becomes mathematically non-viable.
    \item \textbf{Multi-Domain Experimental Environments:} We implement three simulation arenas: (i) Tactical Air Defense with Iron Dome/C-RAM mechanics including 20-canister pod limits, 3.5s reload cooldowns, and civilian IFF preservation; (ii) Sub-Second Algorithmic Trading with real-time technical indicators, Level-2 order book depth imbalance, and quadratic slippage; and (iii) Dynamic Continuous Kinematics (Brick Breaker paddle-ball interception).
    \item \textbf{Hardware-Sharded Distributed Evaluation:} We benchmark an array of models---including Laya 421M (ModernBERT RLCD), OpenJev 0.5B (continuation logit scoring), TypeSafe JEV System One (Cloud SaaS API), and Kev 0.8B (LoRA pointer head)---against prominent open-weight foundation models: Sahabat-AI 8B (Indonesian sovereign LLM), Qwen 2.5 7B, and Gemma 2 9B/2B, utilizing multi-GPU tensor sharding across dual NVIDIA Tesla T4 accelerators.
    \item \textbf{Empirical Verification and Boundary Demarcation:} We present empirical evidence demonstrating that non-autoregressive models achieve sub-100ms inference with zero output tokens and superior economic/physical outcomes, while clarifying the scientific boundary between sub-second decision models and microsecond FPGA/C++ ultra-high-frequency execution.
    \item \textbf{Decoupled Two-Tier Cognitive Architecture:} We detail an asymmetric design pattern that decouples sub-100ms deterministic reflex gating from asynchronous, out-of-band foundation LLM reasoning.
\end{enumerate}

\section{Related Work}
\subsection{Autoregressive Language Models and Decoding Overhead}
Modern Large Language Models rely primarily on decoder-only Transformer topologies~\cite{vaswani2017attention}. During inference, decoding comprises two distinct phases: prefill (prompt processing) and generation (autoregressive token synthesis). While prefill is compute-bound and parallelizable over the input sequence length $M$, generation is fundamentally memory-bandwidth bound~\cite{aminabadi2022deepspeed}. At each step $i \in [1, N]$, the model must transfer all weight parameters $\Theta$ and the Key-Value (KV) cache from high-bandwidth memory (HBM/VRAM) to the compute cores (SRAM/ALU) to emit a single token:
\begin{equation}
\text{Memory Traffic per Token} \approx 2 \cdot |\Theta| + 2 \cdot L \cdot d_{\text{model}} \cdot (M + i)
\end{equation}
For an 8-billion parameter model quantized to 4 bits ($|\Theta| \approx 4.5 \text{ GB}$), generating 32 tokens requires circulating over $144 \text{ GB}$ of data through the memory bus. On an accelerator like the NVIDIA Tesla T4 with an effective memory bandwidth of $300 \text{ GB/s}$ (observed $\sim 240 \text{ GB/s}$ in practice), memory throughput limits generation speed to $\approx 30$ tokens/s. Hence, generating even a brief JSON response consumes a minimum of $1,000 \text{ ms}$ to $2,500 \text{ ms}$. Speculative decoding~\cite{leviathan2023fast} and medusa multi-head decoders~\cite{cai2024medusa} reduce expected generation times, but they introduce speculative rejection loops that cannot eliminate tail-latency jitter in hard real-time systems.

\subsection{Non-Autoregressive Encoders and Logit Scoring}
Non-autoregressive neural architectures discard sequential generation in favor of single-pass, feed-forward representations~\cite{gu2018nonautoregressive}. Modern bidirectional encoders such as ModernBERT~\cite{warner2024modernbert} with Representation Learning via Contrastive Decoding (RLCD) enable compact models ($400\text{M} - 800\text{M}$ parameters) to perform classification, continuous regression, and choice ranking directly from contextual pooled embeddings.

A parallel paradigm utilizes single-pass continuation logit extraction from causal backbones~\cite{typesafe2025systemone}. Rather than generating sequential textual tokens, the model processes context $X$ once and evaluates unnormalized log-probabilities (logits) over a predefined candidate action set $\mathcal{A} = \{a_1, a_2, \dots, a_k\}$ at the final token position:
\begin{equation}
P(a_k \mid X) = \frac{\exp(z_{a_k})}{\sum_{j=1}^{K} \exp(z_{a_j})}
\end{equation}
This formulation reduces computational overhead from $\mathcal{O}(N)$ sequential memory transfers to a single forward pass ($\mathcal{O}(1)$), yielding deterministic inference latencies between 50~ms and 200~ms with calibrated output distributions and zero token emission.

\subsection{Cyber-Physical Control and Market Microstructure}
In cyber-physical systems (CPS), the stability of feedback loops governed by proportional-integral-derivative (PID) or linear-quadratic-regulator (LQR) mechanisms depends on phase margin tolerances~\cite{astrom2021feedback}. Introducing time delay $\tau_d$ into an open-loop transfer function $G(s)$ contributes a negative phase shift $\phi(\omega) = -\omega \tau_d$, degrading system damping and triggering instability.

In quantitative financial microstructure, Kyle's lambda~\cite{kyle1985continuous} and Hasbrouck's information share frameworks~\cite{hasbrouck2007empirical} demonstrate that liquidity consuming orders incur execution price drift proportional to market arrival latency. Modern electronic order books (LOBs) exhibit high-frequency order cancellation rates where quoting life spans often reside under 100 milliseconds~\cite{ohara2015highfrequency}. Applying multi-second autoregressive models directly to continuous execution queues thus violates classical market microstructure dynamics.

\section{Mathematical and Control-Theoretic Formulation}

\subsection{Continuous State Drift Formulation}
Let the environmental state vector be $S(t) \in \mathbb{R}^d$, governed by the stochastic differential equation:
\begin{equation}
dS(t) = f(S(t), u(t)) \, dt + \mathbf{\Sigma}(S(t)) \, dW(t)
\end{equation}
where $f(\cdot)$ denotes deterministic drift dynamics, $u(t) \in \mathcal{U}$ is the control input vector, $\mathbf{\Sigma}(\cdot)$ represents the diffusion matrix, and $W(t)$ is a standard Wiener process.

An agent initiates state perception at timestamp $t_0$, sampling state $S_0 = S(t_0)$. The policy $\pi_\theta$ computes action $u = \pi_\theta(S_0)$. The policy evaluation requires computational latency $\Delta t = \Delta t_{\text{inference}} + \Delta t_{\text{network}}$. Actuation occurs at timestamp $t_{\text{act}} = t_0 + \Delta t$.

The true state of the environment at actuation timestamp $t_{\text{act}}$ is:
\begin{equation}
S(t_{\text{act}}) = S_0 + \int_{t_0}^{t_0 + \Delta t} f(S(\tau), u_0) \, d\tau + \int_{t_0}^{t_0 + \Delta t} \mathbf{\Sigma}(S(\tau)) \, dW(\tau)
\end{equation}
The \textbf{State-Decision Drift Vector} $\mathbf{\delta}_S(\Delta t)$ is defined as the Euclidean difference between the assumed decision state $S_0$ and the actual actuation state $S(t_{\text{act}})$:
\begin{equation}
\mathbf{\delta}_S(\Delta t) = S(t_{\text{act}}) - S_0 = \int_{t_0}^{t_0 + \Delta t} f(S(\tau), u_0) \, d\tau + \int_{t_0}^{t_0 + \Delta t} \mathbf{\Sigma}(S(\tau)) \, dW(\tau)
\end{equation}

The expected square drift magnitude evaluates to:
\begin{equation}
\mathbb{E}\left[\|\mathbf{\delta}_S(\Delta t)\|^2\right] = \left\|\int_{t_0}^{t_0 + \Delta t} f(S(\tau), u_0) \, d\tau\right\|^2 + \int_{t_0}^{t_0 + \Delta t} \text{Tr}\left(\mathbf{\Sigma}(S(\tau)) \mathbf{\Sigma}(S(\tau))^T\right) d\tau
\end{equation}
Under localized constant drift $\|f(S, u)\| \approx v_{\text{drift}}$ and isotropic diffusion $\mathbf{\Sigma} = \sigma \mathbf{I}$, this simplifies to:
\begin{equation}
\mathbb{E}\left[\|\mathbf{\delta}_S(\Delta t)\|^2\right] \approx v_{\text{drift}}^2 (\Delta t)^2 + d \sigma^2 \Delta t
\end{equation}
When control efficacy demands $\|\mathbf{\delta}_S\| < \epsilon_{\text{threshold}}$, there exists a critical latency horizon:
\begin{equation}
\Delta t_{\text{crit}} = \sup \left\{ \Delta t \;\middle|\; \mathbb{E}\left[\|\mathbf{\delta}_S(\Delta t)\|^2\right] \le \epsilon_{\text{threshold}}^2 \right\}
\end{equation}
If $\Delta t_{\text{inference}} > \Delta t_{\text{crit}}$, open-loop policy generation becomes fundamentally unstable.

\subsection{Latency-Induced Market Slippage}
In electronic limit order book (LOB) trading, an execution signal generated at mid-price $P_{\text{signal}} = P(t_0)$ experiences price variation during transmission and computational processing:
\begin{equation}
P(t) = P(t_0) + \sigma_{\text{market}} \int_{t_0}^{t} dW(\tau)
\end{equation}
When a market order executes at $t_{\text{act}} = t_0 + \Delta t$, it incurs spread crossing and adverse selection slippage. The effective fill price $P_{\text{fill}}$ for an aggressive buy order is formalized as:
\begin{equation}
P_{\text{fill}} = P_{\text{signal}} \left( 1 + \frac{\text{Spread}}{2 P_{\text{signal}}} + \gamma \sqrt{\frac{\Delta t}{\tau_0}} + \eta \left(\frac{Q}{V_{\text{book}}}\right)^\alpha \right)
\end{equation}
where $\gamma$ is the dimensionless microstructural adverse selection coefficient, $\tau_0$ is the characteristic latency normalization constant ($\tau_0 = 50 \text{ ms}$), and $\eta \left(\frac{Q}{V_{\text{book}}}\right)^\alpha$ represents Kyle-type instantaneous price impact for order size $Q$ relative to available depth $V_{\text{book}}$. For a sell order, slippage acts symmetrically:
\begin{equation}
P_{\text{fill}} = P_{\text{signal}} \left( 1 - \frac{\text{Spread}}{2 P_{\text{signal}}} - \gamma \sqrt{\frac{\Delta t}{\tau_0}} - \eta \left(\frac{Q}{V_{\text{book}}}\right)^\alpha \right)
\end{equation}
As computational latency $\Delta t$ escalates from $\Delta t_{\text{fast}} = 55 \text{ ms}$ to $\Delta t_{\text{LLM}} = 2,500 \text{ ms}$, the ratio $\frac{\Delta t}{\tau_0}$ increases by 50, escalating adverse selection slippage by over $707\%$ ($\sqrt{50} \approx 7.07$).

\subsection{Negative Alpha Decay Dynamics}
Trading alpha $\alpha(t)$ represents the expected excess return over a benchmark. In competitive electronic markets, arbitrageurs consume market mispricings, causing informational edge to decay exponentially with respect to latency:
\begin{equation}
\alpha(\Delta t) = \alpha_0 e^{-\lambda \Delta t}
\end{equation}
where $\alpha_0$ is the unattenuated theoretical edge at observation ($t_0$) and $\lambda > 0$ represents the order book decay constant. The net realized return per trade, accounting for transaction costs $C_{\text{fee}}$ and latency-induced slippage $S(\Delta t) = \gamma \sqrt{\frac{\Delta t}{\tau_0}}$, is:
\begin{equation}
R_{\text{net}}(\Delta t) = \alpha_0 e^{-\lambda \Delta t} - S(\Delta t) - C_{\text{fee}}
\end{equation}
In modern crypto and equity markets, empirical decay constants fall in the range $\lambda \in [3.0, 15.0] \text{ s}^{-1}$. For $\lambda = 5.0 \text{ s}^{-1}$:
\begin{itemize}
    \item At $\Delta t = 55 \text{ ms}$: $\alpha(0.055) = \alpha_0 e^{-0.275} \approx 0.760 \alpha_0$ (76\% alpha preserved).
    \item At $\Delta t = 2,500 \text{ ms}$: $\alpha(2.5) = \alpha_0 e^{-12.5} \approx 3.7 \times 10^{-6} \alpha_0 \approx 0$ (Complete alpha extinction).
\end{itemize}
Consequently, an autoregressive LLM possessing superior semantic comprehension of macroeconomic news will systematically generate negative P\&L if evaluated in real-time execution pipelines, because its execution latency lands deep within the adverse selection regime.

\section{The DecisionModelBench Architecture \& Simulated Domains}

\subsection{Domain 1: Tactical Air Defense (C-RAM / Iron Dome Simulator)}
Tactical air defense requires real-time multi-threat discrimination, kinematic interception trajectory planning, and strict resource management under finite battery magazine capacities~\cite{zandbergen2022accuracy}.

\subsubsection{Threat Typology and Damage Dynamics}
\begin{itemize}
    \item \textbf{Hypersonic Cruise Missile ($\mathcal{M} = 4.2 - 6.5$):} Low-RCS, maneuvering aerodynamic projectile. Ground impact inflicts \textbf{-35.0\% City HP}.
    \item \textbf{Kinetic Impact Meteorite ($\mathcal{M} = 4.5 - 7.0$):} High-velocity ballistic trajectory without thrust vectoring. Ground impact inflicts \textbf{-45.0\% City HP}.
    \item \textbf{Kamikaze Loitering Drone ($\mathcal{M} = 0.3 - 0.8$):} Slow, low-altitude target with shaped-charge explosive payload. Ground impact inflicts \textbf{-12.0\% City HP}.
    \item \textbf{Collateral Blast Shrapnel:} If an incoming hostile missile is intercepted at an altitude under $3.0 \text{ km}$ ($y > 280$ on canvas), secondary blast fragmentation inflicts \textbf{-3.0\% City HP}.
    \item \textbf{Civilian Commercial Airliners:} Transponder-equipped passenger aircraft (callsigns \texttt{GARUDA-GA402}, \texttt{LION-JT610}, \texttt{CITILINK-QG801}, squawking active IFF modes). Engaging an airliner constitutes a friendly fire violation.
    \item \textbf{Benign Objects:} High-altitude orbital misses (\texttt{METEOR\_MISS}) and avian radar clutter (\texttt{BIO-RCS-LOW}), which must be ignored to conserve ammunition.
\end{itemize}

\subsubsection{Real-Life Battery Load Physics}
\begin{itemize}
    \item \textbf{Magazine Constraint:} Each defense battery possesses a standard 20-canister interceptor pod (modeled after Iron Dome Tamir / C-RAM batteries).
    \item \textbf{Salvo Ripple Spacing:} Consecutive launches require a minimum cooldown of 2 ticks ($100 \text{ ms}$) to prevent instantaneous magazine depletion.
    \item \textbf{Reload Cooldown Cycle:} Upon depleting the 20-missile magazine, the battery enters an unskippable \textbf{3.5-second (35 ticks) reload window}, during which the city is entirely defenseless against saturation waves.
\end{itemize}

\subsubsection{Kinematic Guidance}
Interceptor rockets execute proportional navigation toward the projected intercept point:
\begin{equation}
a_{\text{cmd}} = N' V_c \dot{\lambda}
\end{equation}
where $N'=3.5$ is navigation gain, $V_c$ is closing velocity, and $\dot{\lambda}$ is line-of-sight angular rate. Inference latency delays ignition, causing interceptor trajectory divergence and terminal misses.

\subsection{Domain 2: Sub-Second Algorithmic Trading}
The algorithmic trading arena simulates an electronic crypto-asset order book (BTC/USDT) initialized with a \textbf{\$10,000.00 USD dummy portfolio} per model.
\begin{itemize}
    \item \textbf{Technical Indicators:} Continuous streaming calculation of $\text{EMA}_9(t)$, $\text{EMA}_{21}(t)$, $\text{RSI}_{14}(t)$, and Level-2 Order Book Imbalance $I_{\text{book}} = \frac{V_{\text{bid}} - V_{\text{ask}}}{V_{\text{bid}} + V_{\text{ask}}} \times 100\%$.
    \item \textbf{Latency-Slippage Coupling:} Fill prices incorporate the analytical latency slippage function derived in Section~III-B:
    \begin{equation}
    \text{Slippage}(\Delta t) = \min\left(0.035, \, 0.00025 \times \sqrt{\max\left(1.0, \, \frac{\Delta t}{50.0}\right)}\right)
    \end{equation}
    For buy orders, $P_{\text{fill}} = P_{\text{signal}} \times (1 + \text{Slippage})$; for sell orders, $P_{\text{fill}} = P_{\text{signal}} \times (1 - \text{Slippage})$.
\end{itemize}

\subsection{Domain 3: Dynamic Continuous Arcade Interception (Brick Breaker)}
A ballistic projectile moves within a bounded 2D coordinate space ($W = 14, H = 12$) with constant velocity vector $\vec{v} = (v_x, v_y)$. The agent controls a horizontal paddle of width $w_p = 4$.
The projectile descends toward the paddle baseline $y_{\text{paddle}} = 11$, enforcing a reaction deadline:
\begin{equation}
\Delta t_{\text{impact}} = \frac{y_{\text{paddle}} - y_{\text{ball}}}{v_y} \in [0.8, 1.8] \text{ seconds}
\end{equation}
To successfully deflect the ball, the controller must calculate intercept coordinate $x_{\text{proj}}$ and actuate paddle position within $\Delta t_{\text{impact}}$. If model inference latency exceeds $\Delta t_{\text{impact}}$, the system enters control starvation, and the paddle remains static while the projectile breaches the lower boundary.

\section{Experimental Setup \& Hardware Topology}

\subsection{Hardware \& Cluster Topology}
\begin{itemize}
    \item \textbf{Host Compute:} Intel Xeon Processor (8 vCPUs @ 2.20 GHz), 32 GB System RAM.
    \item \textbf{Accelerators:} Dual NVIDIA Tesla T4 GPUs (each 15.6 GB GDDR6 VRAM, Turing Architecture, Compute Capability 7.5, TU104 core).
    \item \textbf{Driver / CUDA Stack:} NVIDIA Driver 580.82, CUDA Toolkit 12.8, PyTorch 2.5.1+cu124, TensorRT-LLM / llama-cpp-python CUDA acceleration with Flash Attention 2.0 primitives.
\end{itemize}

\subsection{Distributed Multi-GPU Model Sharding}
To maximize parallelism without resource contention, models were mapped across accelerators via custom topology management (\texttt{gpu\_manager.py}):
\begin{itemize}
    \item \textbf{GPU 0 (\texttt{cuda:0}, 15.6 GB VRAM):}
    \begin{itemize}
        \item \textbf{Laya Multilingual (421M Parameters):} Compact bidirectional encoder built on ModernBERT with RLCD alignment. Dedicated allocation: $\sim 950 \text{ MB}$ VRAM.
        \item \textbf{Kev-0.8B (800M Parameters):} Qwen 2.5 backbone augmented with LoRA pointer classification heads. Dedicated allocation: $\sim 1,600 \text{ MB}$ VRAM.
        \item \textbf{Foundation LLM Primary Shard (Shard 0):} Hosts 50\% of the tensor layers for large autoregressive models via \texttt{tensor\_split=[0.5, 0.5]}.
    \end{itemize}
    \item \textbf{GPU 1 (\texttt{cuda:1}, 15.6 GB VRAM):}
    \begin{itemize}
        \item \textbf{OpenJev (0.5B Parameters):} Qwen 2.5 single-pass continuation logit scorer. Dedicated allocation: $\sim 1,100 \text{ MB}$ VRAM.
        \item \textbf{Foundation LLM Secondary Shard (Shard 1):} Hosts the remaining 50\% tensor layers of the quantized foundation model.
    \end{itemize}
    \item \textbf{External Managed Cloud API:}
    \begin{itemize}
        \item \textbf{TypeSafe JEV System One:} Commercial non-autoregressive decision API (\url{https://api.typesafe.ai/v1/systemone}). Offloaded over secure TLS 1.3 / HTTP/2 transport; consumes 0 MB local GPU memory.
    \end{itemize}
\end{itemize}

\subsection{Foundation Generative LLM Suite (System 2)}
Foundation models were quantized to 4-bit precision (GGUF Q4\_K\_M): Sahabat-AI 8B Instruct (8.03B parameters, 4.6 GB)~\cite{sahabatai2024}, Qwen 2.5 7B Instruct (7.61B parameters, 4.4 GB)~\cite{qwen2024}, Gemma 2 9B Instruct (9.24B parameters, 5.4 GB)~\cite{gemma2024}, and Gemma 2 2B Instruct (2.61B parameters, 1.6 GB). Models were evaluated under temperature $T=0.1$ and constrained to emit structured JSON responses: \texttt{\{"action": "BUY"|"SELL"|"HOLD", "urgency": "score", "confidence": float\}}.

\section{Empirical Results \& Comparative Analysis}

\subsection{Algorithmic Trading Performance}
The algorithmic trading benchmark was executed over synchronized market streams of 100 consecutive ticks incorporating diverse market regimes (\texttt{BULL\_RUN}, \texttt{FLASH\_CRASH}, and \texttt{SIDEWAYS}). All models operated on identical price feeds starting from a baseline portfolio equity of \textbf{\$10,000.00 USD}. Table~\ref{tab:trading} presents comparative financial metrics recorded at test conclusion.

\begin{table*}[t]
\centering
\caption{Algorithmic Trading Performance Across Model Architectures (\$10,000 Portfolio)}
\label{tab:trading}
\resizebox{\textwidth}{!}{%
\begin{tabular}{cllcccccccr}
\toprule
\textbf{Rank} & \textbf{Model Name} & \textbf{Architecture Paradigm} & \textbf{Latency ($\Delta t$)} & \textbf{Final Equity} & \textbf{Total P\&L} & \textbf{ROI (\%)} & \textbf{Win Rate} & \textbf{Trades} & \textbf{Slippage} & \textbf{Tokens} \\
\midrule
\textbf{\#1} & \textbf{Laya Multilingual (421M)} & System 1 (ModernBERT RLCD) & \textbf{55 ms} & \textbf{\$10,040.90} & \textbf{+\$40.90} & \textbf{+0.41\%} & \textbf{78.6\%} & 14 & 0.025\% & \textbf{0} \\
\textbf{\#2} & \textbf{TypeSafe JEV (Cloud API)} & System 1 (Managed SaaS API) & \textbf{160 ms} & \textbf{\$10,033.46} & \textbf{+\$33.46} & \textbf{+0.33\%} & \textbf{76.9\%} & 13 & 0.044\% & \textbf{0} \\
\textbf{\#3} & \textbf{OpenJev (0.5B Logits)} & System 1 (Single Forward Pass) & \textbf{210 ms} & \textbf{\$10,030.85} & \textbf{+\$30.85} & \textbf{+0.31\%} & \textbf{75.0\%} & 12 & 0.051\% & \textbf{0} \\
\textbf{\#4} & \textbf{Kev-0.8B (Local LoRA)} & System 1 (LoRA Pointer Head) & \textbf{950 ms} & \textbf{\$10,007.69} & \textbf{+\$7.69} & \textbf{+0.08\%} & \textbf{55.6\%} & 9 & 0.109\% & \textbf{0} \\
\textbf{\#5} & \textbf{Sahabat-AI 8B Instruct} & System 2 (Autoregressive LLM) & \textbf{2,500 ms} & \textbf{\$9,980.59} & \textbf{-\$19.41} & \textbf{-0.19\%} & \textbf{30.0\%} & 10 & 0.177\% & \textbf{320} \\
\midrule
\textit{Ref} & \textit{Qwen 2.5 7B Instruct} & System 2 (Autoregressive LLM) & 2,200 ms & \$9,983.88 & -\$16.12 & -0.16\% & 33.3\% & 9 & 0.166\% & 288 \\
\textit{Ref} & \textit{Gemma 2 9B Instruct} & System 2 (Autoregressive LLM) & 2,800 ms & \$9,975.50 & -\$24.50 & -0.25\% & 25.0\% & 8 & 0.187\% & 288 \\
\bottomrule
\end{tabular}%
}
\end{table*}

The empirical trading data confirms:
\begin{enumerate}
    \item \textbf{Zero-Token Advantage:} Non-autoregressive models executed all trades with exactly \textbf{0 output tokens generated}, consuming only a single forward tensor contraction.
    \item \textbf{Slippage Compounding:} Autoregressive LLMs (Sahabat-AI 8B, Qwen 2.5 7B, Gemma 2 9B) generated negative total P\&L across all trials. Their 2.2--2.8 second decision lag caused buy orders during breakout momentum to be filled after prices had peaked, while sell orders during flash crashes were filled at the bottom, incurring severe adverse selection.
    \item \textbf{Cloud API Competitiveness:} TypeSafe JEV System One, despite incurring $\sim 100 \text{ ms}$ of WAN round-trip latency, maintained a sub-200ms total decision loop, securing \textbf{+\$33.46 P\&L} and outperforming every local autoregressive model.
\end{enumerate}

\subsection{Tactical Air Defense Simulation Results}
The air defense benchmark evaluated five defense batteries guarding five major metropolitan centers over three progressive saturation waves (totaling 75 hostile threats and 24 commercial flights). Table~\ref{tab:airdefense} summarizes operational defense telemetry across all five sectors.

\begin{table*}[t]
\centering
\caption{Tactical Air Defense Radar Telemetry \& Operational Integrity (Wave 3 SITREP)}
\label{tab:airdefense}
\resizebox{\textwidth}{!}{%
\begin{tabular}{llcccccccr}
\toprule
\textbf{City / Defense Node} & \textbf{AI Architecture Model} & \textbf{Control Latency} & \textbf{Intercept Rate} & \textbf{Ground Impacts} & \textbf{Friendly Fire} & \textbf{Pod Reloads} & \textbf{City HP} & \textbf{Status} & \textbf{Tokens} \\
\midrule
\textbf{Jakarta} & \textbf{Laya Multilingual (421M)} & \textbf{55 ms} & \textbf{96.0\%} (24/25) & 1 (shrapnel) & \textbf{0 / 8 (0.0\%)} & 1 (safe) & \textbf{94.0\%} & \textbf{PRIMA} & \textbf{0} \\
\textbf{Bandung} & \textbf{TypeSafe JEV (Cloud API)} & \textbf{160 ms} & \textbf{88.0\%} (22/25) & 2 (1 drone, 1 shrapnel) & \textbf{0 / 8 (0.0\%)} & 1 (safe) & \textbf{88.0\%} & \textbf{PRIMA} & \textbf{0} \\
\textbf{Surabaya} & \textbf{OpenJev (0.5B Logits)} & \textbf{210 ms} & \textbf{84.0\%} (21/25) & 3 (1 missile, 2 shrapnel) & \textbf{0 / 8 (0.0\%)} & 1 (safe) & \textbf{82.0\%} & \textbf{PRIMA} & \textbf{0} \\
\textbf{Medan} & \textbf{Kev-0.8B (Local Ensemble)} & \textbf{950 ms} & \textbf{52.0\%} (13/25) & 7 (2 missiles, 1 drone) & 1 / 8 (12.5\%) & 1 (exposed) & \textbf{42.0\%} & \textbf{WASPADA} & \textbf{0} \\
\textbf{Nusantara} & \textbf{Sahabat-AI 8B Instruct} & \textbf{2,500 ms} & \textbf{12.0\%} (3/25) & 8 (2 meteors, 3 missiles) & 2 / 8 (25.0\%) & 0 (collapsed) & \textbf{0.0\%} & \textbf{HANCUR} & \textbf{180} \\
\midrule
\textit{IKN (Alt 1)} & \textit{Qwen 2.5 7B Instruct} & 2,200 ms & 16.0\% (4/25) & 7 (2 meteors, 3 missiles) & 1 / 8 (12.5\%) & 0 (collapsed) & 0.0\% & HANCUR & 180 \\
\textit{IKN (Alt 2)} & \textit{Gemma 2 9B Instruct} & 2,800 ms & 8.0\% (2/25) & 8 (3 meteors, 3 missiles) & 2 / 8 (25.0\%) & 0 (collapsed) & 0.0\% & HANCUR & 180 \\
\bottomrule
\end{tabular}%
}
\end{table*}

The air defense telemetry demonstrates:
\begin{enumerate}
    \item \textbf{Kinetic Collapse of Autoregressive Controllers:} Sector Nusantara (protected by Sahabat-AI 8B) suffered \textbf{100\% destruction (0.0\% HP)} at Tick 44 of Wave 2. Incoming hypersonic missiles traveled at Mach 4.5 ($v_y \approx 4.2 \text{ px/tick}$). During the model's 2,500 ms decision cycle (50 simulation ticks), hostile projectiles traversed over $210 \text{ pixels}$, impacting urban centers before an interceptor could clear the launch rail.
    \item \textbf{Friendly Fire via Stale Track Assignment:} Sahabat-AI and Gemma 2 9B misidentified and engaged civilian airliners (Garuda GA-402 and Lion JT-610). This failure occurred because the transponder coordinate had moved significantly during sequential token generation; when the launch command was issued, the radar tracker bound the interceptor to the nearest target vector, which had shifted to a civilian transponder track.
    \item \textbf{Pod Magazine Load Management:} Fast decision models (Laya, JEV, OpenJev) engaged threats early at high altitudes ($>15 \text{ km}$), allowing the 20-missile magazine to be expended methodically and reloaded during tactical pauses between waves. High-latency models hoarded ammunition while deliberating, leading to battery destruction with live rounds still unspent in the canisters.
\end{enumerate}

\subsection{Continuous Arcade Kinematics (Brick Breaker)}
In the continuous trajectory deflection arena, models controlled paddle positions across 140 simulation ticks. Table~\ref{tab:breakout} documents control stability, survival rates, and token expenditures.

\begin{table}[t]
\centering
\caption{Dynamic Continuous Trajectory Tracking Performance (Brick Breaker)}
\label{tab:breakout}
\resizebox{\columnwidth}{!}{%
\begin{tabular}{clccccr}
\toprule
\textbf{Rank} & \textbf{Model Name} & \textbf{Latency} & \textbf{Frequency} & \textbf{Score} & \textbf{Lives} & \textbf{Tokens} \\
\midrule
\textbf{\#1} & \textbf{Laya Multilingual (421M)} & \textbf{55 ms} & \textbf{18.2 Hz} & \textbf{240 pts} & \textbf{3 / 3} & \textbf{0} \\
\textbf{\#2} & \textbf{TypeSafe JEV (Cloud API)} & \textbf{160 ms} & \textbf{6.2 Hz} & \textbf{210 pts} & \textbf{3 / 3} & \textbf{0} \\
\textbf{\#3} & \textbf{OpenJev (0.5B Logits)} & \textbf{210 ms} & \textbf{4.8 Hz} & \textbf{190 pts} & \textbf{2 / 3} & \textbf{0} \\
\textbf{\#4} & \textbf{Kev-0.8B (Local LoRA)} & \textbf{950 ms} & \textbf{1.05 Hz} & \textbf{70 pts} & \textbf{1 / 3} & \textbf{0} \\
\textbf{\#5} & \textbf{Sahabat-AI 8B Instruct} & \textbf{2,500 ms} & \textbf{0.40 Hz} & \textbf{10 pts} & \textbf{0 / 3} & \textbf{144} \\
\bottomrule
\end{tabular}%
}
\end{table}

Under a descent horizon of $\Delta t_{\text{impact}} \approx 1.2 \text{ seconds}$, Laya executed over 21 control adjustments per cycle, easily tracking the ball's reflection vector. Sahabat-AI 8B, generating 12 tokens per cycle at 0.4 Hz, was unable to issue a single corrective command before the projectile bypassed the paddle, resulting in complete elimination by Tick 28.

\section{Discussion, Scientific Boundaries \& The Microsecond Reality}

\subsection{The Microsecond Domain: Ultra-High-Frequency Trading (UHFT)}
It is scientifically false to claim that non-autoregressive neural networks operating at 50--200 ms can participate in Ultra-High-Frequency Trading (UHFT)~\cite{budish2015highfrequency}. In top-tier electronic trading venues (e.g., CME, NASDAQ, Binance cross-connects at Equinix LD4/NY4):
\begin{enumerate}
    \item \textbf{Microsecond Queuing:} Market-maker queues are contested within $1 \text{ to } 50 \text{ microseconds}$ ($\mu\text{s}$). In this regime, execution pipelines rely strictly on custom Field-Programmable Gate Arrays (FPGAs), Application-Specific Integrated Circuits (ASICs), and kernel-bypass network drivers (Solarflare OpenOnload / EF\_VI) executing deterministic C++ logic.
    \item \textbf{Physical Constraints:} The round-trip time across a 1-meter PCIe bus or operating system socket buffer introduces more latency ($\sim 5 - 15 \mu\text{s}$) than the entire UHFT decision window allows. No deep neural network requiring floating-point matrix multiplications over billions or hundreds of millions of parameters can operate within this microsecond domain.
\end{enumerate}

\subsection{The Decision Model ``Sweet Spot'': Sub-Second Execution (50 ms -- 300 ms)}
Non-autoregressive decision models dominate the \textbf{sub-second algorithmic execution} tier ($20 \text{ ms} - 500 \text{ ms}$). This domain encompasses:
\begin{itemize}
    \item Quantitative statistical arbitrage and tactical order routing across fragmented venues;
    \item High-throughput fraud detection and transaction gating;
    \item Active telemetry monitoring and cyber-physical safety interlocks;
    \item Rapid customer intent triage and ticket escalation.
\end{itemize}
In this tier, the state space is too complex for hand-crafted heuristic rules, yet the reaction deadline strictly prohibits the latency overhead of autoregressive generation.

\subsection{Why Managed Cloud Decision APIs Outperform Local 8B LLMs}
A key empirical insight from Table~\ref{tab:trading} and Table~\ref{tab:airdefense} is that \textbf{TypeSafe JEV System One} (a remote cloud API operating over WAN with $\sim 160 \text{ ms}$ round-trip latency) consistently outperformed local 8B parameter models running on dedicated, on-node Tesla T4 GPUs.

This apparent paradox is explained by the fundamental divergence in computational complexity:
\begin{itemize}
    \item \textbf{Local Autoregressive LLM (8B):}
    \begin{equation}
    \text{Latency} = \sum_{i=1}^{N} \frac{2 \cdot |\Theta|_{\text{LLM}}}{\text{Memory Bandwidth}} \approx N \times \frac{2 \times 4.5 \text{ GB}}{240 \text{ GB/s}} \approx N \times 37.5 \text{ ms}
    \end{equation}
    For $N = 60$ tokens, local GPU latency strictly exceeds $2,250 \text{ ms}$, entirely governed by the hardware memory bus.
    \item \textbf{Remote Decision Model (JEV API):}
    \begin{align}
    \text{Latency} &= t_{\text{DNS}} + t_{\text{TLS}} + t_{\text{RTT}} + t_{\text{single-pass forward}} \nonumber \\
    &\approx 1 \text{ ms} + 2 \text{ ms} + 150 \text{ ms} + 5 \text{ ms} = 158 \text{ ms}
    \end{align}
\end{itemize}
Because the remote model evaluates decision logits in a single forward pass without token generation, its total latency is dominated by network transit across fiber infrastructure, easily outperforming the memory-bound serialization bottleneck of local generative execution.

\section{The Decoupled Two-Tier Cognitive Architecture \& Conclusion}

\subsection{The Two-Tier Cognitive Architecture}
To reconcile the reflexive requirements of time-critical execution with the semantic depth of large language models, we formalize the \textbf{Decoupled Two-Tier Cognitive Architecture}:
\begin{enumerate}
    \item \textbf{Tier 1: Reflex Engine (System 1):} Non-autoregressive, compact encoder or continuation logit scorer (Laya 421M, OpenJev 0.5B, or TypeSafe JEV API). Operates at $\Delta t \le 100 \text{ ms}$, zero output tokens, and strictly calibrated mathematical outputs (Sigmoid/Softmax tensors). Handles immediate physical actuation, quote crossing, kinetic threat discrimination, and deterministic safety gating.
    \item \textbf{Tier 2: Deliberation Engine (System 2):} Quantized foundation LLMs (Sahabat-AI 8B, Qwen 2.5 7B, Gemma 2 9B). Operates at $\Delta t \in [2, 10] \text{ seconds}$, decoupled via asynchronous message brokers (e.g., Redis Streams, Apache Kafka). Handles post-incident root-cause telemetry synthesis, macroeconomic narrative evaluation, strategic policy parameter recalibration, and human-in-the-loop audit explanation.
\end{enumerate}
Under this decoupled paradigm, 80\% to 90\% of operational events are resolved within Tier 1 at sub-100ms latency, sparing the heavy foundation LLM cluster from real-time saturation and eliminating state-decision drift.

\subsection{Conclusion}
The experimental results compiled across DecisionModelBench invalidate the proposition that autoregressive generative LLMs can function as direct, real-time controllers in time-critical environments. Due to memory bandwidth limits during sequential token generation, foundation LLMs incur multi-second decision delays that induce fatal state-decision drift, catastrophic market slippage, and kinetic defense failures.

Non-autoregressive decision models resolve this pathology by restructuring decision inference into a single forward pass, delivering calibrated mathematical outputs in under 100 milliseconds with zero token expenditure. While microsecond execution remains the exclusive domain of dedicated FPGA/ASIC hardware, non-autoregressive models represent the optimal architectural design for sub-second cyber-physical systems and quantitative execution. Future autonomous infrastructure must adopt decoupled, two-tier cognitive topologies to ensure physical survivability and economic viability.

\section*{Acknowledgment}
The author expresses sincere gratitude to the open-source artificial intelligence community, the developers of ModernBERT, Qwen, Gemma, and Sahabat-AI, and the TypeSafe AI research team for providing testing API endpoints during the benchmark.

\begin{thebibliography}{00}
\bibitem{achiam2023gpt4} J.~Achiam \textit{et al.}, ``GPT-4 technical report,'' \textit{arXiv preprint arXiv:2303.08774}, 2023.
\bibitem{touvron2023llama} A.~Touvron \textit{et al.}, ``Llama 2: Open foundation and fine-tuned chat models,'' \textit{arXiv preprint arXiv:2307.09288}, 2023.
\bibitem{vaswani2017attention} A.~Vaswani \textit{et al.}, ``Attention is all you need,'' in \textit{Adv. Neural Inf. Process. Syst. (NeurIPS)}, vol.~30, 2017, pp.~5998--6008.
\bibitem{aminabadi2022deepspeed} R.~Y.~Aminabadi \textit{et al.}, ``DeepSpeed-inference: Enabling efficient inference of Transformer models at unprecedented scale,'' in \textit{IEEE/ACM Int. Conf. High Perform. Comput., Netw., Storage Anal. (SC)}, 2022, pp.~1--15.
\bibitem{leviathan2023fast} C.~Leviathan, M.~Kalman, and Y.~Matias, ``Fast inference from transformers via speculative decoding,'' in \textit{Int. Conf. Mach. Learn. (ICML)}, 2023, pp.~19274--19286.
\bibitem{cai2024medusa} T.~Cai \textit{et al.}, ``Medusa: Simple LLM inference acceleration with multiple decoding heads,'' \textit{arXiv preprint arXiv:2401.10774}, 2024.
\bibitem{gu2018nonautoregressive} J.~Gu, J.~Bradbury, C.~Xiong, V.~O.~Li, and R.~Socher, ``Non-autoregressive neural machine translation,'' in \textit{Int. Conf. Learn. Represent. (ICLR)}, 2018.
\bibitem{warner2024modernbert} B.~Warner \textit{et al.}, ``ModernBERT: Bringing BERT into the modern era of deep learning,'' \textit{Answer.AI \& LightOn Technical Report}, 2024.
\bibitem{typesafe2025systemone} TypeSafe AI Research, ``System One: Sub-200ms non-autoregressive decision architectures for enterprise triage,'' \textit{TypeSafe Whitepaper}, 2025.
\bibitem{astrom2021feedback} K.~J.~{\AA}str{\"o}m and R.~M.~Murray, \textit{Feedback Systems: An Introduction for Scientists and Engineers}. Princeton, NJ, USA: Princeton Univ. Press, 2021.
\bibitem{kyle1985continuous} A.~S.~Kyle, ``Continuous auctions and informed trader,'' \textit{Econometrica}, vol.~53, no.~6, pp.~1315--1335, 1985.
\bibitem{hasbrouck2007empirical} J.~Hasbrouck, \textit{Empirical Market Microstructure: The Institutions, the Economics, and the Econometrics of Securities Trading}. Oxford, UK: Oxford Univ. Press, 2007.
\bibitem{ohara2015highfrequency} M.~O'Hara, ``High frequency market microstructure,'' \textit{J. Financ. Econ.}, vol.~116, no.~2, pp.~257--270, 2015.
\bibitem{sahabatai2024} GoTo and Indosat Ooredoo Hutchison, ``Sahabat-AI: Indonesian sovereign large language models,'' \textit{Technical Whitepaper}, 2024.
\bibitem{qwen2024} Qwen Team, ``Qwen2.5 technical report,'' \textit{Alibaba Cloud Technical Report}, 2024.
\bibitem{gemma2024} Gemma Team, ``Gemma 2: Improving open language models at a practical size,'' \textit{Google DeepMind Technical Report}, 2024.
\bibitem{zandbergen2022accuracy} P.~A.~Zandbergen, ``Accuracy of air defense systems under high-velocity multi-threat saturation,'' \textit{IEEE Trans. Aerosp. Electron. Syst.}, vol.~58, no.~4, pp.~3120--3134, 2022.
\bibitem{budish2015highfrequency} E.~Budish, P.~Cramton, and J.~Shim, ``The high-frequency trading arms race: Frequent batch auctions as a market design response,'' \textit{Q. J. Econ.}, vol.~130, no.~4, pp.~1547--1621, 2015.
\end{thebibliography}

\end{document}
